% Standalone by default. In an article with amsmath and amsthm loaded,
% define \RDecorationFragment before inputting this file.
\ifdefined\RDecorationFragment\else
\documentclass[11pt]{article}
\usepackage{amsmath,amsthm}
\fi
\newtheorem{rdtheorem}{Theorem}
\newtheorem{rdlemma}{Lemma}
\newcommand{\rdSL}{\mathbf{SL}}
\newcommand{\rdLG}{\mathbf{LG}}
\newcommand{\rdLI}{\mathbf{LI}}
\newcommand{\rdtyp}{\mathrm{typ}}
\newcommand{\rdcnt}{\operatorname{alph}}
\ifdefined\RDecorationFragment\else
\begin{document}
\title{State decoration and short distribution representatives}
\author{R16 research note}\date{23 November 2026}\maketitle
\fi

Write $c(w)=\rdcnt(w)$, and $b_j(w)$ for the whole word with a short tag if
$|w|<\max(1,2(j-1))$, otherwise for its tagged prefix and suffix of length
$j-1$. Let $N_j$ count all length-$j$ factors with overlaps. A typing is a
finite positive-word semigroup morphism, and $s\preceq h$ means that equality
of $h$ implies equality of $s$. An admissible $a$-context has $a$ holes in
order, each used once, positive internal separators and possibly empty outer
pieces. Tuples have positive components. Equivalence means equality both
of component types and of the sets of accepting admissible contexts. A target
is $(f,s)$-substitutable if, in every arity $1,\ldots,f$, equally typed tuples
sharing one accepting context are equivalent. A good grammar has positive
template components, linear nondeleting, nonpermuting, nonmerging rules:
each child's variables occur once, in order, and consecutive variables of
that child cannot be adjacent in the same output component. Reduced means
reachable and productive. Its original thickness $\theta\ge1$ is the maximum
of one and, over rules, the minimum length of a successful sentence using
that rule. Write $v$ for the number of nonterminals and $r$ for the rank cap.
The quantity $\tau(G,s)$ takes the maximum, over realized typed symbols
other than the fresh start of the refinement,
of the shortest length in the corresponding equivalence class; the representative
need not be generated there. All generated tuples of a fixed type at one
symbol share a derivational context: sibling components are positive and
good rules preserve the order and separation of holes. Empty target words
are handled separately.

We use two prior results with the following interfaces.
The corrected local-counting theorem (R14, Theorem L) says that if positive
membership depends on $(b_m,N_m)$ and the target is $(f,t)$-substitutable
with $t\preceq(c,b_m)$, then on a reduced good presentation over alphabet $B$,
\[
 \tau=O_{B,m,f,r}(v^{D+2}\theta^{D+1}),\qquad D=f|B|^m .
\]
The local-group short-yield theorem (compendium, theorem on local groups and corollary on type-level shortest yields) supplies a \emph{generated} tuple of length
$C v\theta$ for each realizable vector of values of a fixed morphism into
a finite local group, at its original nonterminal. A finite local group is
a semigroup whose $eTe$ is a group for every idempotent $e$.
The results below use these established statements as inputs.

\begin{rdtheorem}[Transfer]
Let $s$ be fixed, and let a reduced good fan-out-$\le f$, rank-$\le r$
grammar present the positive part of an $(f,s)$-substitutable language $L$.
Let a fixed deterministic complete machine have states $Q$, initial $q_0$
and positive transition semigroup $T$. Its trace $D_p(w)$ emits $(q,a)$
on reading $a$ in state $q$, so $\pi(q,a)=a$ preserves length.
Assume positive membership depends on $(\rho(w),N_\ell(w))$, where
$\rho$ is a fixed finite semigroup morphism and
$s\preceq(c,b_\ell,\rho)$. Suppose for some fixed $d$ the input state,
output state, $c(D_p(w))$ and $b_d(D_p(w))$ determine $\rho(w)$ on
all realizable positive traces with arbitrary reachable entry $p$.
More generally one may replace the preceding membership, typing and recovery
assumptions by these two operational hypotheses at the width $m$ below:
(i) membership in the initial-state trace language is determined by
$(b_m,N_m)$; (ii) on coherent components $t=(c,b_m)$ determines the
actual projected $s$-type. For a one-state machine this version also works
with $m=1$ (coherence and endpoint matching are then automatic).
For each realizable componentwise vector of entry/exit pairs at an original
nonterminal assume a generated witness there of total length at most $\beta$.
Set $m=\max(2,d,\ell)$, $D=f(|Q||\Sigma|)^m$.
Then, with constants depending only on fixed parameters,
\[
 \tau(G,s)=O\big(v^{D+2}[v(\theta+r\beta)]^{D+1}\big).
\]
If $T\in\rdLG$, this is $O(v^{3D+4}\theta^{D+1})$.
\end{rdtheorem}
\begin{proof}
Drop unreachable states. Put $L^\#=D_{q_0}(L\cap\Sigma^+)$ and
$t=(c,b_m)$ over $\Gamma=Q\times\Sigma$. A word is a legal global trace
iff its first entry is $q_0$ and all adjacent letters match the transition.
The data $(b_m,N_m)$ determine the first and last letters and all counts
down to length one: to go down one order count prefixes of the longer
windows and add the last shorter window, treating short words literally.
This tests legality and supplies $c,b_d$ and both endpoints.
Recovery gives $\rho$ and letterwise projection of the length-$\ell$ counts
gives the desired input counts. Hence $L^\#$ is locally counting.
On every coherent component $t$ alone gives the endpoints, its input
$c,b_\ell$ and, by recovery, $\rho$, hence its actual input $s$-type. These are exactly the two facts supplied
directly in the operational version. Equal component endpoints follow from
$b_m$ for $m\ge2$, and are automatic for one state.

Fix any arity $1\le a\le f$, componentwise equal-$t$ positive tuples $X,Y$,
a shared accepting decorated context $C_0$, and any decorated test $C_1$.
The four words are $C_0[X],C_0[Y],C_1[X],C_1[Y]$. The first two give
internal coherence. With $x=\pi X,y=\pi Y$ the four input words are
\[ (\pi C_0)[x],\quad(\pi C_0)[y],\quad(\pi C_1)[x],\quad(\pi C_1)[y]. \]
Internal gaps stay positive. The input types match componentwise, hence
membership of the last two input words agrees. A failure inside a fixed
piece of $C_1$, at its start, or at a boundary is simultaneous since
component endpoints match, even with empty outside pieces. Otherwise the
two outputs are both coherent from $q_0$. Their acceptance then agrees
by projection. If no $C_0$ exists the premise of substitutability is empty.
This proves all decorated rectangles.

Build the grammar over $\Gamma$ with symbols $(A;(p_i,q_i)_i)$.
For every original rule choose entry/exit vectors for its parent and
children. Scan each output template from its entry: emit tagged terminals,
require a child's stated entry to equal the current state at its variable,
then resume at its exit, and check the terminal state of the template.
Use a unary fresh start joining initial $q_0$ to each possible final state.
Induction on trees proves equality with all decorated generated tuples
having the specified endpoints. Conversely erasing states gives an actual
original derivation. Trim. A productive derivation at a reachable symbol
can be inserted at that symbol, so trimming loses no needed witness.
Rank and fan-out caps are preserved, as are goodness and positivity.
There are $v'\le1+v|Q|^{2f}$ symbols, and the grammar size is
$O_{Q,f,r}(1+\|G\|)$: each original rule has a bounded number of choices,
each with the original terminal length. This includes the unary start rule.

A rule writes at most $\theta$ terminals, since a successful original sentence
using it contains them all. Each surviving decorated symbol has a generated
tuple of length at most $\beta$ by choosing the short endpoint witness and
tracing it from its specified entries (allow $\max(1,\beta)$ for notation).
The fresh start also has such a witness whenever the target is nonempty.
To witness a particular surviving rule, reach its parent by a simple path
in the reachable dependency graph, choose the actual rules along it and use
short generated tuples for all siblings and final children. There are at
most $v'+1$ rules on this spine. These substitutions use the same decorated
nonterminals, and do not appeal to language substitutability. Therefore
$\bar\theta_{G^\#}\le(v'+1)(\theta+r\beta)$ (for an empty grammar the
claim is vacuous). Theorem L applies with alphabet $\Gamma$ and the above $t$.

Finally take a successful occurrence of any realized original typed symbol,
with tuple $x$. The sentence trace gives a generated tuple $X$ at an actual
decorated symbol and a common accepting derivational context $F$ for its
realized $t$-sort. Theorem L supplies a positive tuple $Y$ of the same
$t$ and decorated distribution, within its length bound. In particular
$F[X],F[Y]$ are accepted. Thus $Y$ is coherent, and $x,\pi Y$ have equal
actual component $s$-types. For each original test $E$, the four input
words are
\[ (\pi F)[x],\quad(\pi F)[\pi Y],\quad E[x],\quad E[\pi Y]. \] The shared first
pair and original substitutability identify membership of the second pair.
This does not require the same decorations to work in $E$. Length is
unchanged by $\pi$. Substitution of the numerical bounds proves the first
claim. A realized endpoint vector has a witnessing tuple of full transition
values. Preserving these values preserves every specified endpoint, so
the local-group theorem gives $\beta=O(v\theta)$ when $T\in\rdLG$.
\end{proof}

\paragraph{A sharper accounting.}
The proof of Theorem L uses its thickness variable only to bound explicit
terminals per rule and the length of one untyped generated tuple per symbol.
Its fixed outer context is used for acceptance, never charged to the final
length. Indeed height and the number of marked ancestors depend only on
the number of types and symbols; each DFS increment charges terminals
written on one ancestor path, and each impure-window colour charges the
short tuples at the cut. The skeleton charges explicit terminals.
Thus one may substitute $b=\max(1,\theta,\beta)$ for the thickness variable
in that proof, while still obtaining the full rule-thickness estimate above
when it is needed. This gives $O(v^{D+2}b^{D+1})$, hence
$O(v^{2D+3}\theta^{D+1})$ for a local-group lower machine.

\section*{Which regular skeletons admit recovery?}
Products below act on the right. Write $U\wr_Q T$ for pairs $(F,t)$,
$F:Q\to U$, with multiplication
$(F,t)(H,s)=(q\mapsto F(q)H(qt),ts)$.
The positive image is used throughout, even when $Q=T^1$ contains a
formally adjoined initial state. The symbol $*$ denotes the pseudovariety
generated by these cascades, with the right factor reading input first.
Finite products of cascades embed in a single cascade using the product
lower action and tuple of upper functions. Division means a quotient of a
subsemigroup. Standard wreath-product division or the displayed formula
allows the regular lower action on $T^1$: for an arbitrary action, replace
$F_a(t)$ by the vector $(F_a(qt))_{q\in Q}$ and evaluate at $t=1$ to
recover the original functions. The upper semilattice remains a semilattice.

\begin{rdlemma}
$\rdLG*\rdLG=\rdLG$. In particular upper group counters or a locally trivial
boundary register above a local-group machine still give an LG positive image.
\end{rdlemma}
\begin{proof}
Take an idempotent $E=(F,b)$ in a wreath product with both factors in $\rdLG$
and an idempotent $Z=(H,z)$ in its local monoid. Then $z=b$ since
$z$ is an idempotent of $bTb$. At every $u$ fixed by $b$, both $F(u)$
and $H(u)$ are idempotents, and $EZ=ZE=Z$ puts $H(u)$ in
$F(u)UF(u)$. Hence $H(u)=F(u)$. At arbitrary $u$, $ub$ is fixed, so
$H(u)=F(u)H(ub)=F(u)F(ub)=F(u)$. A finite monoid with just one idempotent
is a group (a power of every element is the identity).
This proves the inclusion; the reverse uses a trivial upper factor.
\end{proof}

\begin{rdtheorem}[Exact skeleton class for existential full recovery]
A regular positive-word skeleton $R$ admits an LG lower machine and a
fixed width giving the component recovery in the transfer theorem for
its entire syntactic morphism $\rho$, if and only if its syntactic
semigroup belongs to $\rdSL*\rdLG$.
\end{rdtheorem}
\begin{proof}
For necessity consider the trace from the initial state. Its content
and its boundary type can be accumulated by a finite semilattice and
a positive boundary semigroup, respectively. The latter is locally trivial:
an idempotent has a long tag and sandwiching any positive tag between two
copies leaves it unchanged. Endpoints are obtained from the lower machine.
Consequently $R$ is recognized by the product of a semilattice-over-LG
cascade and a locally-trivial-over-LG cascade. Empty accumulations can be
handled by an initial state, without changing the positive acting semigroup.
The lemma and division by the syntactic congruence prove necessity.

Conversely lift the syntactic morphism letter by letter into a divisor
presentation in the regular wreath product $U\wr_{T^1}T$, with $U$ a
finite semilattice, $T\in\rdLG$. Thus each letter $a$ has a chosen lift
$(F_a,t_a)$ in the covering subsemigroup. Put $k=|T|$. We use the
local-group structure statement that $T^k$ is its completely simple
minimal ideal $K$. In particular for $p\in K,x\in T^1$ there exists
$h\in T^1$ with $hxp=p$. Indeed $xp$ and $p$ have the same column in a
Rees matrix representation of $K$: writing $p=e p$ for a left identity
$e\in K$ shows $xp=(xe)p$. In a fixed column one can left-multiply from
any row and group coordinate to any other by the Rees multiplication law.

Run the lower state $x\in T^1$ by $x\mapsto xt_a$ and, for each event
$(x,a)$, maintain its own counter modulo $k+1$, incremented just after that
event. This is a finite abelian group above $T$, hence an LG lower machine
by the lemma. It emits its complete state and the input letter.
We now decode the \emph{whole wreath element} for a component, from
\emph{any} entrance into this machine, using content and boundary width
$k+1$ (enlarge to at least two). For each event its observed pre-increment
residues visit consecutive residues around its cycle. Their number of
distinct values gives its total count truncated at $k+1$, irrespective of
the initial counters. These values are available from the decorated content.
If the component has fewer than $k$ letters read it directly.
Otherwise read its first $k$ transitions and let $p$ be their input product
in $T$. From each truncated count subtract the corresponding prefix count
and test positivity. This gives exactly the set of events occurring
\emph{after} the prefix, even when a count saturates, since at most $k$
occurrences were discarded. No order-sensitive upper products are changed.

Let $x$ now denote the component's actual initial $T^1$ state. Fix
$h=h(x,p)$ with $hxp=p$. For each possible initial coordinate $y\in T^1$
in the original wreath representation, simulate the first $k$ letters
starting at $y$, multiplying their upper values. At each subsequent event
$(z,a)$ in the observed tail the correct simulated pre-state is $yhz$:
indeed $z=xpt$ for the product $t$ of the intervening input.
Multiply in $F_a(yhz)$ once per distinct tail event. This reconstructs the
upper value at $y$ because the upper semigroup is commutative and idempotent.
There is always an upper factor for a positive word, so empty subproducts
cause no difficulty. The lower value of the whole input is $hz_{\rm end}$
from the observed exit $z_{\rm end}=xpt_{\rm tail}$. All coordinates of
the lift, hence its syntactic image, are recovered. This is one decoder on
all entrances with fixed finite memory, independent of the grammar and
count filter. Sufficiency follows.
\end{proof}

\paragraph{Regular threshold counting.}
In positive semigroup notation,
\[
 \rdSL*\rdLG=\mathbf{LT}*\rdLG=\mathbf{LTT}*\rdLG .
\]
Indeed, Straubing (1985), Theorem 4.3 and Section 7.1, gives
$\mathbf{LT}=\rdSL*\mathbf D$; taking the upper monoid variety to
be finite commutative aperiodic monoids in Theorem 4.3 gives
$\mathbf{LTT}=\mathbf{ACom}*\mathbf D$.
For the latter specialization each generator's powers stabilize at a
common threshold; conversely all truncated factor counts are recorded
by a product of finite truncated-addition monoids. End windows and
short-word tags are included. The semigroup reducts of
$\mathbf{ACom}$ belong to $\rdSL*\mathbf G$: on a fixed generating
alphabet keep each letter count modulo its stabilization threshold
(or a larger common modulus $n\ge1$). From the content of letters
decorated by these counters, one obtains $\min(n,|w|_a)$ by counting
distinct pre-residues for $a$. This recognizes every fibre from the
identity, and hence the morphism divides such a cascade. As
$\mathbf D,\mathbf G\subseteq\rdLG$, associativity of transformation
wreath products and the lemma give the equality, using
$\rdSL\subseteq\mathbf{LT}\subseteq\mathbf{LTT}$ for the reverse
inclusions. Exact unbounded factor counts in the target are separate.

The inclusion $\mathbf{LTT}*\mathbf G\subseteq\rdSL*\rdLG$ is strict.
Take the language with an even number of contiguous factors $ab$.
Its positive syntactic semigroup is in $\rdLG$: keep the previous letter
in a right-zero lower semigroup and toggle an upper group bit on $ab$.
In any purported finite aperiodic-over-group recognition, choose $N$ a
positive multiple of the lower group exponent. A lift of the block
$B=a^N b^N$ has the form $(F,1)$, so its powers stabilize by aperiodicity
at each of the finitely many upper coordinates. Yet membership of $B^j$
alternates with $j$. This also excludes $\mathbf{LT}*\mathbf G$.

\paragraph{References to the algebraic interfaces.}
H. Straubing, \emph{Finite semigroup varieties of the form $V*D$},
J. Pure Appl. Algebra 36 (1985), 53--94, Theorems 3.1 and 4.3 and
Section 7.1. J.-E. Pin, \emph{Mathematical Foundations of Automata
Theory}, transformation wreath products and the wreath product principle,
Theorem 4.6, pp. 291--292 in the consulted version.
\ifdefined\RDecorationFragment\else\expandafter\enddocument\fi
